We thank the reviewer for the careful reading and thoughtful questions.

\medskip
\noindent\textbf{Q1 (Technical novelty).}\quad
We apologize for not being precise enough in the coupling argument presentation: one intermediate step was missing in the draft wording.

More precisely, there is no technical novelty for the existence/uniqueness proof following Nguyen et al.
The only technical novelty comes from modifying the coupling template to our channel-wise sum (their three-layer assumption / coupling display): after introducing the coupling $\pi_t$ and using bounded first-layer activation, one obtains a bound of the form
$$
\mathbb{E}_{\pi_t}\left[\left|\left\langle \mathbb{E}_{Z}\left[ \Delta_{2}^{H*}(t,\cdot)-\Delta_{2}^{H*}(\cdot;\bar w)\mid X=x \right],\varphi_{1}(\langle u,x\rangle) \right\rangle\right|\right]
\le K\,\mathbb{E}_{\pi_t}\left[\left|\Delta_{2}^{H*}(t,\cdot)-\Delta_{2}^{H*}(\cdot;\bar w)\right|\right],
$$
and one then controls the right-hand side by weighted coupled gaps (same logic as Nguyen \& Pham), with low-rank-specific channel aggregation.

Our low-rank translation of this missing step is that the main coupling term is controlled through a channel sum under $\pi_t$:
$$
\mathbb{E}_{\pi_t}\Bigl[(1+|\bar w_2|)\,|\bar w_2|\,\sum_{k=1}^{r} |\bar w_{1,k}|\,|w_{1,k}-\bar w_{1,k}|\Bigr]\to 0,
$$
with notation matched to our neuronal embedding variables in Eq.~(1).
This is exactly the place where the RF-LR structure replaces the fully connected contraction chain by channelwise control plus mixing bounds.

We will make this explicit in the appendix by adding:
(i)~the post-\emph{translates-to} display;
(ii)~the time-regularity condition for each channel,
$$
\operatorname*{ess\,sup}_{t}\left|\partial_t w_1(t,U_k)\right|\to 0,\qquad k=1,\dots,r;
$$
(iii)~the bi-Lipschitz/bounded-mixing constants used to close the backward estimate.

For ReLU, the closure step uses the sharp inequality
$$
\left|\operatorname{ReLU}\left(\sum_{i=1}^{r}x_i\right)-\operatorname{ReLU}\left(\sum_{i=1}^{r}y_i\right)\right|
\le \left|\sum_{i=1}^{r}(x_i-y_i)\right|
\le \sum_{i=1}^{r}|x_i-y_i|
$$
rather than any extra factor-of-$r$ slack.
We will state this explicitly (with compact-support / boundedness assumptions for pre-activations) to clarify why $r$-channel coupling closes.

\medskip
\noindent\textbf{Q2 (Spectral bias and Section~4).}\quad
On universality and training regimes: we will revise Section~4 to separate statements cleanly.
In the NTK / lazy training regime, the effective bias is governed by the kernel spectrum and one expects the classical low-frequency-first behavior emphasized in that line of work.
The empirical ``tilt'' toward comparatively higher-frequency structure that we highlight is observed in a richer feature-learning setting for our low-rank RF architecture; we do not claim that this behavior is universal across all scalings (including finite width outside mean-field parameterizations or lazy limits).

Regarding Rahaman et al., ``On the Spectral Bias of Neural Networks'' (2019), arXiv:1806.08734, \url{https://arxiv.org/abs/1806.08734}, is an apt reference for the standard Fourier / low-frequency-first picture in sufficiently wide, kernel-like regimes.

On why low-rank maps might carry more high-frequency content (strictly heuristic, work in progress / under review):
in depth-$L$ ReLU networks, one can view NN outputs through ``sum over paths'' (per-region linear parametrizations); shrinking the effective rank reduces how many independent hinge directions can co-activate across depth, which loosely relaxes constraints on how linear regions can tile input space---informally, a less constrained CPWL can allow relatively more high-dimensional faces versus low-dimensional ones and thus alter how energy is distributed across spatial frequencies in a Fourier lens.
This complements the shorter ``fewer effectively independent hinge directions / boundary geometry'' phrasing we will keep in the main text and will be marked explicitly as exploratory reasoning and ongoing review, not as a formal implication.

On exposition and scope: we agree Section~4 is currently qualitative and that the two-point analytic anchor is minimal.
We will tone down claims, acknowledge the simplicity of the two-point construction, and add citations for the CPWL / compositional picture and Fourier heuristics beyond Rahaman et al.
